\input{../../stock/en-tete_v6.tex}
\newcommand{\titrechapitre}{Erreurs d'arrondi}
\newcommand{\numerochapitre}{1}

\renewcommand{\headrulewidth}{0.4pt}
\renewcommand{\footrulewidth}{0.4pt}
\fancyfoot[L]{Notes de Raphaël JONTEF}
\fancyfoot[C]{\thepage}
\fancyfoot[R]{Enseirb-Matmeca}
\fancyhead[L]{Cours d'algorithmique numérique}
\fancyhead[R]{\titrechapitre}
        \usepackage{hhline}

\begin{document}
\input{../../stock/commands.tex}

% Page de titre custom
\setlength{\headheight}{15pt}
\begin{titlepage}
    \centering
    \vspace*{3cm}
    {\huge Chapitre \numerochapitre\par}
    {\Huge \bfseries\titrechapitre\par}
    \vspace{2cm}
    {\Large Notes rédigées par Raphaël JONTEF\par}
{\Large Avec l'aide de Github Copilot\par}
    \vspace{0.5cm}
    {\normalsize Cours enseigné par Marc Durufle\par}
    \vspace{4cm}
    {\small Enseirb-Matmeca\par}
    {\small filière informatique\par}
    \vfill
    {\large \today\par}
\end{titlepage}

% plan du cours
\tableofcontents
\newpage


% -----------Corps du document--------------


\section{Résolution de systèmes linéaires}

Soit $A \in \matrice{n}(\R)$, $B \in \R^n$
En pratique, la solution numérique $x + \delta x$ de l'égalité vectorielle $AX = B$ vérifie~:
$$A\times(x+\delta x) = b + \delta b$$
On a alors~:
$$\delta b = A\times(x + \delta x) - b$$
Le vecteur $x$ étant inconnu, on ne peut pas calculer directement $\delta b$. On cherche donc à majorer $\norme{\delta b}$.\\\par

\subsection{Rappels sur les normes usuelles sur $\R^n$}

\begin{definition}{}{norme infinie}
    Pour $X \in \R^n$, on définit la norme infinie par~:
    $$\norme{X}_\infty = \max_{1 \leq i \leq n} |x_i|$$
    En Python, on peut la calculer avec la fonction \texttt{numpy.linalg.norm(X, inf)}.
\end{definition}

\begin{definition}{}{norme de Mannhattan}
    Pour $X \in \R^n$, on définit la norme de Manhattan par~:
    $$\norme{X}_1 = \sum_{i=1}^n |x_i|$$
    En Python, on peut la calculer avec la fonction \texttt{numpy.linalg.norm(X, 1)}.
\end{definition}

\begin{definition}{}{norme euclidienne}
    Pour $X \in \R^n$, on définit la norme euclidienne par~:
    $$\norme{X}_2 = \left(\sum_{i=1}^n |x_i|^2\right)^{1/2}$$
    En Python, on peut la calculer avec la fonction \texttt{numpy.linalg.norm(X)}.
\end{definition}

\begin{definition}{}{norme subordonnée}
    Pour $A \in \matrice{n}(\R)$, on définit la norme subordonnée associée à une norme vectorielle $\norme{\cdot}$ par~:
    $$\norme{A} = \sup_{X \in \R^n, X \neq 0} \frac{\norme{AX}}{\norme{X}}$$
    En Python, on peut la calculer avec la fonction \texttt{numpy.linalg.norm(A, ord)} où \texttt{ord} est l'ordre de la norme vectorielle associée (\code{1}, \code{2} ou \code{inf}).
    
\end{definition}


\begin{remarque}{}{sur la norme euclidienne d'une matrice}
    On a $\norme{A}_2 = \sqrt{\rho(AA^\top)}$ où $\rho(M)$ est la valeur propre de $M$ de plus grande valeur absolue.
    Cette norme est difficile à calculer en pratique. On préfère donc utiliser les normes $\norme{\cdot}_1$ ou $\norme{\cdot}_\infty$. qui sont plus rapides à calculer.
    
\end{remarque}

\subsection{Erreur d'approximation et méthode d'estimation de cette erreur}

Si $A$ est inversible, on a~:
\begin{align*}
    b &= Ax
    &\Rightarrow \norme{b} &= \norme{Ax} \leq \norme{A}\norme{x}\\
\end{align*}

De même~:
\begin{align*}
    A \times (x + \delta x) &= b + \delta b \\
    &\Rightarrow A \delta x = \delta b\\
    &\Rightarrow \norme{\delta b} = \norme{A \delta x} \leq \norme{A}\norme{\delta x}\\
    &\Rightarrow \norme{b} \norme{\delta b} \leq \norme{A}\norme{x} \norme{A^{-1}}\norme{\delta x}\\
    &\Rightarrow \underbrace{\frac{\norme{\delta x}}{\norme{x}}}_{\text{erreur relative sur la solution}} \leq \underbrace{\norme{A}\norme{A^{-1}}}_{\mr{cond}(A)} \underbrace{\frac{\norme{\delta b}}{\norme{b}}}_{\text{erreur que l'on mesure}}\\
\end{align*}
L'erreur est amplifiée d'un facteur appelé \notion{conditionnement} de la matrice $A$, noté~:
$$\mr{cond}(A) = \norme{A}\norme{A^{-1}}$$
En Python, on le calcule avec \code{cond(A, ord)} (\textit{cf.} plus haut).

\begin{definition}{}{mauvais conditionnement}
    On dit d'une matrice $A \in \matrice{n}(\R)$ qu'elle est \notion{mal conditionnée} si son conditionnement est élevé~:
    \begin{itemize}
        \item Si $\mr{cond}(A) = 10^4$, alors $A$ est mal conditionnée
        \item Si $\mr{cond}(A) = 3$, aors $A$ est bien conditionnée.
    \end{itemize}
\end{definition}

Moralement, pour résoudre $Ax = b$~:
\begin{itemize}
    \item On trouve $x + \delta x$
    \item On évalue le \notion{résidu} $\norme{\delta b} = \norme{b - A\times(x + \delta x)}$
\end{itemize}

\begin{remarque}{}{mauvais conditionnement}
    Les matrices obtenues en discrétisant des équations physiques sont généralement mal conditionnées.
\end{remarque}



\subsection{Calculs de complexité}
Les additions, soustractions et multiplication ont aujourd'hui des coûts temporels semblables, mais la division est plus complexe. On cherche à s'en passer le plus possible.

\begin{exemple}{}{complexité du produit d'une matrice avec un vecteur}
    Soit $A \in \matrice{n,p}(\R)$ et $x \in \R^n$
    Notons $y = Ax$ et $y = (y_1,\, \dots,\, y_n)^\top$. On a~:
    $$\forall i \in \intint{1}{n},\, y_i = \sum_{j=1}^p a_{i,j} x_j$$.
    Le calcul de chaque $y_i$ demande donc $2p$ opérations, ainsi le calcul de $y$ demande $2np$ opérations.
    
\end{exemple}

\begin{exemple}{}{résolution d'un système linéaire triangulaire}
    Soit $T \in \matrice{n,p}(\R)$ et $b \in \R^n$. Calculons la complexité de la résolution de $Tx = b$. On prend $T$ triangulaire inférieure. Ainsi~:
    \begin{align*}
        Tx = b &\Leftrightarrow \begin{cases}
            t_{1,1}x_1 = b_1\\
            t_{2,1}x_1 + t_{2,2}x_2 = b_2\\
            \vdots\\
            t_{n,1}x_1 + t_{n,2}x_2 + \dots + t_{n,n}x_n = b_n
        \end{cases}\\
        &\Leftrightarrow \begin{cases}
            x_1 = \frac{b_1}{t_{1,1}}\\
            x_2 = \frac{b_2 - t_{2,1}x_1}{t_{2,2}}\\
            \vdots\\
            x_n = \frac{b_n - \sum_{j=1}^{n-1} t_{n,j}x_j}{t_{n,n}}
        \end{cases}
    \end{align*}
    Les coûts opératoires pour chaque ligne vérifient le tableau suivant~:
    \begin{center}
        \begin{tabular}{c|c|c|c|c|}
            \textbf{Ligne} & \textbf{additions} & \textbf{soustractions} &\textbf{multiplications} & \textbf{divisions}\\
            \hline
            $1$ & $0$ & $0$ & $0$ & $1$\\
            \hline
            $2$ & $0$ & $1$ & $1$ & $1$\\
            \hline
            $3$ & $0$ & $2$ & $2$ & $1$\\
            \hline
            \vdots & \vdots & \vdots & \vdots & \vdots\\
            \hline
            $n$ & $0$ & $n-1$ & $n-1$ & $1$\\
            \hhline{=|=|=|=|=|}
            \textbf{Somme} & $0$ & $\displaystyle\sum_{i=0}^{n-1} i$ & $\displaystyle\sum_{i=0}^{n-1} i$ & $n$\\
        \end{tabular}
    \end{center}
    On a donc $n^2$ opérations au total.
    
    
\end{exemple}

\subsection{Méthodes de résolution de $Ax = b$}

\subsubsection{Par factorisation de $A$}

\begin{exemple}{}{produit de matrices triangulaires}
    Si $A = LU$  avec $L$ et $U$ triangulaires, alors pour résoudre $Ax = b$, on résout successivement~:
    \begin{itemize}
        \item $Ly = b$ (coût $n^2$)
        \item $Ux = y$ (coût $n^2$)
    \end{itemize}
    Le coût total est donc de $2n^2$ opérations.
\end{exemple}

\subsection{Par méthode itérative}
On cherche à construire une suite $(x^{(k)})_{k\in\N}$ telle que~:
$$\lim_{k \to +\infty} x^{(k)} = x$$
où $x$ est la solution de $Ax = b$. On arrête le processus itératif lorsque le \notion{résidu relatif} $\frac{\norme{\delta b }}{\norme{b}}= \frac{\norme{b - A x^{(k)}}}{\norme{b}}$ est inférieur à un seuil fixé $\varepsilon$~:
$$\frac{\norme{x-x_k}}{\norme{x}} = \frac{\norme{\delta x}}{\norme{x}} \leq \mr{Cond}(A)\epsilon$$

\begin{exemple}{}{de choix d'$\varepsilon$}
    Si on veut une erreur relative sur la solution inférieure à $10^{-6}$ et que le conditionnement de $A$ est de $10^4$, on choisit~:
    $$\varepsilon = \frac{10^{-6}}{10^4} = 10^{-10}$$
\end{exemple}

\begin{remarque}{}{avantages de cette méthode}
    Cette méthode présente plusieurs avantages~:
    \begin{itemize}
        \item pas de nécessité de stocker la matrice $A$ en mémoire
        \item pas de nécessité de calculer une factorisation de $A$ qui pourrait présenter un phénomène de \notion{remplissage} (beaucoup de zéros dans $A$ deviennent des non-zéros dans la factorisation)
    \end{itemize}
\end{remarque}

\begin{remarque}{}{désavantages de cette méthode}
    Cette méthode présente aussi des inconvénients~:
    \begin{itemize}
        \item le nombre d'itérations nécessaires pour atteindre la précision souhaitée peut être élevé et il augmente avec le conditionnement de $A$.
        \item le choix de la méthode itérative et de son paramétrage (ex: préconditionneur) peut être délicat.
    \end{itemize}
    Pour éviter ces inconvénients, on peut tenter de précodnitionner la matrice $A$ pour améliorer son conditionnement, ce qui revient à résoudre~:
    $$M^{-1}Ax = M^{-1}b$$
    où $M$ est un \notion{préconditionneur} choisi pour que $M^{-1}A$ soit mieux conditionnée que $A$.
\end{remarque}

\section{Factorisation de Cholensky}

Soit $A \in \matrice{n}(\R)$ une matrice symétrique définie positive. On cherche à factoriser $A$ sous la forme~:
$$A = LL^\top$$
où $L$ est une matrice triangulaire inférieure.\\

\begin{remarque}{}{matrices symétriques définies positives usuelles}
    Une matrice symétrique $A \in S_n(\R)$ à diagonale strictement dominante~:
    $$\forall i \intint{1}{n},\, a_{ii} > \sum_{j \neq i} \abs{a_{ij}}$$
    est définie positive. C'est un bon moyen de construire des matrices définies positives en pratique, ou de vérifier l'appartenance à $S_n^+$.
\end{remarque}

Une telle matrice $L$, on a~:

$$A = \begin{pmatrix}
    l_{11} & 0 & 0 & \cdots & 0\\
    l_{21} & l_{22} & 0 & \cdots & 0\\
    l_{31} & l_{32} & l_{33} & \cdots & 0\\
    \vdots & \vdots & \vdots & \ddots & \vdots\\
    l_{n1} & l_{n2} & l_{n3} & \cdots & l_{nn}
\end{pmatrix}\times \begin{pmatrix}
    l_{11} & l_{21} & l_{31} & \cdots & l_{n1}\\
    0 & l_{22} & l_{32} & \cdots & l_{n2}\\
    0 & 0 & l_{33} & \cdots & l_{n3}\\
    \vdots & \vdots & \vdots & \ddots & \vdots\\
    0 & 0 & 0 & \cdots & l_{nn}
\end{pmatrix}$$

D'où~:
$$A = \begin{pmatrix}
    l_{11}^2 & l_{11}l_{21} & l_{11}l_{31} & \cdots & l_{11}l_{n1}\\
    l_{11}l_{21} & l_{21}^2 + l_{22}^2 & l_{21}l_{31} + l_{22}l_{32} & \cdots & l_{21}l_{n1} + l_{22}l_{n2}\\
    l_{11}l_{31} & l_{21}l_{31} + l_{22}l_{32} & l_{31}^2 + l_{32}^2 + l_{33}^2 & \cdots & l_{31}l_{n1} + l_{32}l_{n2} + l_{33}l_{n3}\\
    \vdots & \vdots & \vdots & \ddots & \vdots\\
    l_{11}l_{n1} & l_{21}l_{n1} + l_{22}l_{n2} & l_{31}l_{n1} + l_{32}l_{n2} + l_{33}l_{n3} & \cdots & \sum_{k=1}^n l_{nk}^2
\end{pmatrix}$$

En comparant les coefficients de la première colonne, on obtient~:
$$\begin{cases*}
    a_{11} = l_{11}^2\\
    a_{21} = l_{11}l_{21}\\
    a_{31} = l_{11}l_{31}\\
    \vdots\\
    a_{n1} = l_{11}l_{n1}
\end{cases*} \Rightarrow \begin{cases*}
    l_{11} = \sqrt{a_{11}}\\
    l_{i1} = \frac{a_{i1}}{l_{11}} \quad \text{pour } i \in \intint{2}{n}
\end{cases*}$$

Puis pour la deuxième colonne~:
$$\begin{cases*}
    a_{22} = l_{21}^2 + l_{22}^2\\
    a_{32} = l_{21}l_{31} + l_{22}l_{32}\\
    \vdots\\
    a_{n2} = l_{21}l_{n1} + l_{22}l_{n2}
\end{cases*} \Rightarrow \begin{cases*}
    l_{22} = \sqrt{a_{22} - l_{21}^2}\\
    l_{i2} = \frac{a_{i2} - l_{21}l_{i1}}{l_{22}} \quad \text{pour } i \in \intint{3}{n}
\end{cases*}$$

Puis pour la $i$-ème colonne~:
$$\begin{cases*}
    a_{ii} = \sum_{k=1}^i l_{ik}^2\\
    a_{(i+1)i} = \sum_{k=1}^i l_{(i+1)k}l_{ik}\\
    \vdots\\
    a_{ni} = \sum_{k=1}^i l_{nk}l_{ik}
\end{cases*} \Rightarrow \begin{cases*}
    l_{ii} = \sqrt{a_{ii} - \sum_{k=1}^{i-1} l_{ik}^2}\\
    l_{ji} = \frac{a_{ji} - \sum_{k=1}^{i-1} l_{jk}l_{ik}}{l_{ii}} \quad \text{pour } j > i
\end{cases*}$$


\begin{exemple}{}{factorisation de Cholensky}
    Tentons de calculer la factorisation de Cholensky de la matrice~:
    $$A = \begin{pmatrix}
        4 & -2 & -4 \\
        -2 & 10 & 5 \\
        -4 & 5 & 6
    \end{pmatrix}$$

    En travaillant dans l'ordre croissant des colonnes, on trouve~:
    $$L = \begin{pmatrix}
        2 & 0 & 0\\
        -1 & 3 & 0\\
        -2 & 1 & 1
    \end{pmatrix}$$
\end{exemple}


\section{Implémentation de $\matrice{n,p}(\R)$}

Dans cette section nous étudions différentes méthodes de stockage des matrices en mémoire.

\subsection{Le plus fréquent : \code{CSR}}

La méthode \notion{Compressed Sparse Row} (\notion{CSR}) consiste à stocker uniquement les éléments non nuls de la matrice, ainsi que leurs indices de ligne et de colonne. On utilise trois tableaux~:
\begin{enumeratebf}
    \item \code{data} : tableau des valeurs non nulles
    \item \code{indices} : tableau des indices de colonnes des valeurs non nulles
    \item \code{ptr} : tableau des indices dans \code{data} où commence chaque ligne
\end{enumeratebf}

\end{document}